\documentclass[11pt,a4paper]{article}

\usepackage[margin=1in]{geometry}
\usepackage[T1]{fontenc}
\usepackage[utf8]{inputenc}
\usepackage{lmodern}
\usepackage{amsmath,amssymb}
\usepackage{booktabs}
\usepackage{enumitem}
\usepackage{setspace}
\usepackage{xcolor}
\usepackage{hyperref}

\hypersetup{
  colorlinks=true,
  linkcolor=blue!60!black,
  urlcolor=blue!60!black
}

\setlength{\parskip}{0.5em}
\setlength{\parindent}{0pt}
\onehalfspacing

\title{ZPOe 25/26L 2026: Filters On Human Faces\\\large Technical Report on Filter Computation Methodologies}
\author{Student: Adar}
\date{\today}

\begin{document}
\maketitle
\tableofcontents
\newpage

\section{Introduction}
This report presents the implementation methodology for the image processing modules used in the project. The emphasis is on algorithmic computation, numerical handling, and per-filter technical behavior.

\section{Project Structure}
The implementation follows a modular architecture:
\begin{itemize}[leftmargin=1.2cm]
  \item \texttt{facefiltering/filters/*.py}: one module per filter with \texttt{apply(...)}
  \item \texttt{facefiltering/ops.py}: shared NumPy operations (convolution, median, morphology, histogram equalization helpers)
  \item \texttt{facefiltering/gray.py}: grayscale and channel-format helpers
  \item \texttt{facefiltering/validate.py}: input and parameter validation
  \item \texttt{facefiltering/registry.py}: dispatch and categorization
\end{itemize}

\section{Shared Numerical Strategy}
All filters use common safeguards:
\begin{itemize}[leftmargin=1.2cm]
  \item Input normalization to contiguous BGR uint8
  \item Parameter clamping and odd kernel enforcement
  \item Border padding for neighborhood operations
  \item NaN/Inf sanitization where required
  \item Final clipping/normalization to \([0,255]\) and uint8 conversion
\end{itemize}

\section{Detailed Per-Filter Computation}

\subsection{Sobel (magnitude)}
\textbf{Purpose:} detect edge strength from first-order derivatives.
\[
G_x = I * S_x,\quad G_y = I * S_y,\quad M = \sqrt{G_x^2 + G_y^2}
\]
Implementation:
\begin{enumerate}[leftmargin=1.2cm]
  \item Convert image to grayscale.
  \item Convolve with Sobel kernels \(S_x,S_y\).
  \item Compute magnitude map with \(\texttt{hypot}\).
  \item Normalize to uint8 range.
\end{enumerate}

\subsection{Laplacian (abs)}
\textbf{Purpose:} detect high curvature / second-derivative intensity changes.
\[
\nabla^2 I = \frac{\partial^2 I}{\partial x^2} + \frac{\partial^2 I}{\partial y^2}
\]
Implementation:
\begin{enumerate}[leftmargin=1.2cm]
  \item Grayscale conversion.
  \item Convolve with Laplacian kernel.
  \item Apply absolute value.
  \item Normalize to uint8.
\end{enumerate}

\subsection{Canny Edge Detection}
\textbf{Purpose:} robust binary edge extraction with edge thinning and connectivity logic.
Pipeline:
\begin{enumerate}[leftmargin=1.2cm]
  \item Gaussian smoothing
  \item Gradient magnitude/direction
  \item Non-maximum suppression
  \item Double-threshold hysteresis (strong/weak edge propagation)
\end{enumerate}
This implementation computes each stage explicitly in NumPy.

\subsection{Gaussian Blur}
\textbf{Purpose:} low-pass smoothing.
\[
G(x,y)=\frac{1}{2\pi\sigma^2}\exp\left(-\frac{x^2+y^2}{2\sigma^2}\right),\quad I' = I * G
\]
Implementation:
\begin{enumerate}[leftmargin=1.2cm]
  \item Build normalized Gaussian kernel from \(\sigma\).
  \item Convolve each channel.
  \item Clip to uint8.
\end{enumerate}

\subsection{Median}
\textbf{Purpose:} suppress impulse noise while preserving edges.
\[
I'(x,y)=\operatorname{median}\{I(u,v)\mid (u,v)\in \mathcal{N}_{k\times k}(x,y)\}
\]
Implementation uses sliding windows and per-window median on each channel.

\subsection{Unsharp Mask}
\textbf{Purpose:} sharpen details by adding high-frequency residual.
\[
I_{\text{sharp}} = I + a\,(I - G_\sigma * I)
\]
Implementation:
\begin{enumerate}[leftmargin=1.2cm]
  \item Compute Gaussian-blurred image.
  \item Compute detail residual \(I-\text{blur}\).
  \item Add scaled detail using factor \(a\).
  \item Clip to uint8.
\end{enumerate}

\subsection{Binary Threshold}
\textbf{Purpose:} intensity-based segmentation.
\[
I'(x)=
\begin{cases}
255, & I(x)\ge T\\
0, & I(x)\lt T
\end{cases}
\]
Implementation is an element-wise comparison and mapping to \(\{0,255\}\).

\subsection{Gamma}
\textbf{Purpose:} nonlinear tone remapping.
\[
I' = 255\left(\frac{I}{255}\right)^{1/\gamma}
\]
Implementation:
\begin{enumerate}[leftmargin=1.2cm]
  \item Precompute 256-entry lookup table.
  \item Apply LUT to each pixel/channel.
\end{enumerate}

\subsection{Histogram Equalization}
\textbf{Purpose:} improve contrast via cumulative-distribution remapping.
Implementation:
\begin{enumerate}[leftmargin=1.2cm]
  \item Compute luminance estimate.
  \item Build histogram and CDF.
  \item Derive equalization LUT.
  \item Remap luminance and reconstruct color intensities.
\end{enumerate}

\subsection{Morphological Dilation}
\textbf{Purpose:} expand bright structures.
\[
(I\oplus B)(x)=\max_{b\in B} I(x-b)
\]
Implementation uses an elliptical structuring element and local maximum selection for each neighborhood.

\subsection{Morphological Erosion}
\textbf{Purpose:} shrink bright structures.
\[
(I\ominus B)(x)=\min_{b\in B} I(x+b)
\]
Implementation mirrors dilation but applies local minimum selection.

\subsection{High-pass (Fourier)}
\textbf{Purpose:} suppress low frequencies and enhance fine detail.
\[
F=\mathcal{F}\{I\},\quad F' = H_{\text{hp}}\cdot F,\quad I'=\mathcal{F}^{-1}\{F'\}
\]
Implementation:
\begin{enumerate}[leftmargin=1.2cm]
  \item FFT of padded grayscale image.
  \item Center spectrum and apply Gaussian high-pass mask.
  \item Inverse FFT, take real part, normalize.
\end{enumerate}

\subsection{Wiener Deconvolution}
\textbf{Purpose:} restoration from blur/noise model.
\[
\hat{F}(u,v)=\frac{H^*(u,v)}{|H(u,v)|^2 + K}\,G(u,v)
\]
Implementation:
\begin{enumerate}[leftmargin=1.2cm]
  \item Build Gaussian PSF from \texttt{psf\_size}.
  \item FFT of image and PSF transfer function.
  \item Compute Wiener gain with noise term \(K\).
  \item Inverse FFT and clip to uint8 output.
\end{enumerate}

\section{Methodology Coverage Summary}
\begin{itemize}[leftmargin=1.2cm]
  \item \textbf{Spatial convolution:} Sobel, Laplacian, Gaussian, Unsharp
  \item \textbf{Spatial non-linear:} Median, Canny
  \item \textbf{Point-wise intensity:} Threshold, Gamma
  \item \textbf{Global remapping:} Histogram Equalization
  \item \textbf{Morphological:} Dilation, Erosion
  \item \textbf{Frequency-domain:} High-pass Fourier, Wiener
\end{itemize}

\section{Conclusion}
The implemented modules form a complete and methodologically diverse classical image processing suite. Each filter is isolated, computationally explicit, and numerically stabilized, enabling both technical reproducibility and clear algorithm-level analysis.

\end{document}
